% --- Archon physics formalization source begin ---
% archon:physics
% archon:covers PhyXMiniProblems/problem_phyx_mini_0338.lean
% archon:source-report reports/phyx_mini/problem_phyx_mini_0338.source.json

\chapter{Physics problem phyx\_mini\_0338}
\label{ch:PhyXMiniProblems_problem_phyx_mini_0338}

\paragraph{Problem source.}
\#\# Physical scenario

The \$pV\$-diagram in \textbackslash{}textbf\{figure\} shows a process \$abc\$ involving \$0.450\$ mol of an ideal gas.

\#\# Question

How much heat had to be added during the process to increase the internal energy of the gas by \$15\{,\}000\textasciitilde{}\textbackslash{}text\{J\}\$?

\#\# Answer choices

- A: A: 3.4 \textbackslash{}times 10\textasciicircum{}5J

- B: B: 3.6 \textbackslash{}times 10\textasciicircum{}4J

- C: C: 2.6 \textbackslash{}times 10\textasciicircum{}4J

- D: D: 3.2 \textbackslash{}times 10\textasciicircum{}4J

\#\# Figure caption (auxiliary; use the image as primary evidence)

The image is a pressure-volume (PV) diagram. Here's a detailed description:

- **Axes**: 
  - The y-axis represents pressure, labeled as \textbackslash{}( p \textbackslash{}) in units of \textbackslash{}( \textbackslash{}text\{Pa\} \textbackslash{}times 10\textasciicircum{}5 \textbackslash{}), with values from 0 to 8.
  - The x-axis represents volume, labeled as \textbackslash{}( V \textbackslash{}), in units of cubic meters (m\textbackslash{}(\textasciicircum{}3\textbackslash{})), with values from 0 to 0.080.

- **Points**:
  - Three points labeled as \textbackslash{}( a \textbackslash{}), \textbackslash{}( b \textbackslash{}), and \textbackslash{}( c \textbackslash{}) are marked on the graph:
    - Point \textbackslash{}( a \textbackslash{}) is approximately at (0.020 m\textbackslash{}(\textasciicircum{}3\textbackslash{}), \textbackslash{}( 2.0 \textbackslash{}times 10\textasciicircum{}5 \textbackslash{}) Pa).
    - Point \textbackslash{}( b \textbackslash{}) is approximately at (0.060 m\textbackslash{}(\textasciicircum{}3\textbackslash{}), \textbackslash{}( 4.0 \textbackslash{}times 10\textasciicircum{}5 \textbackslash{}) Pa).
    - Point \textbackslash{}( c \textbackslash{}) is approximately at (0.060 m\textbackslash{}(\textasciicircum{}3\textbackslash{}), \textbackslash{}( 8.0 \textbackslash{}times 10\textasciicircum{}5 \textbackslash{}) Pa).

- **Lines and Arrows**:
  - A blue line with an arrow leads from point \textbackslash{}( a \textbackslash{}) to point \textbackslash{}( b \textbackslash{}), indicating a change in state with increasing volume and pressure.
  - A second blue line with an arrow leads vertically from point \textbackslash{}( b \textbackslash{}) to point \textbackslash{}( c \textbackslash{}), showing an increase in pressure at a constant volume.

- **Grid**:
  - The background has a grid aiding in reading specific values of pressure and volume more precisely.

The diagram likely represents a process in a thermodynamic cycle, illustrating changes in pressure and volume.

\#\# Dataset metadata

Laws of Thermodynamics; Physical Model Grounding Reasoning, Multi-Formula Reasoning

\paragraph{Recorded answer/context.}
B: B: 3.6 \textbackslash{}times 10\textasciicircum{}4J

\paragraph{Figure/image path.}
/root/proposal\_for\_physic/science-mango/phyx\_mini\_run/phyx\_data/test\_image/338.png

\paragraph{Formalization target.}
create a compiling Lean file with sorry bodies at `PhyXMiniProblems/problem\_phyx\_mini\_0338.lean`. The Lean declarations must preserve the physical quantities, dimensions or dimensional roles, figure labels, governing-law hypotheses, and final relation expressed by this problem.
Use Mathlib/Physlib names found through LeanExplore where available. If a physics API is missing, introduce faithful local abstractions rather than scalar placeholder aliases.

\begin{theorem}[Physics formalization target]
\label{thm:physics:phyx_mini_0338:target}
\lean{PhyXMiniProblems.ProblemPhyXMini0338.recordedAnswerChoice_isCorrect}
\uses{def:physics:phyx-mini-0338:phyxminiproblems-problemphyxmini0338-idealgasprocessabc, def:physics:phyx-mini-0338:phyxminiproblems-problemphyxmini0338-matchesproblemandprimaryfigure, def:physics:phyx-mini-0338:phyxminiproblems-problemphyxmini0338-hasphysicalprocessparameters, def:physics:phyx-mini-0338:phyxminiproblems-problemphyxmini0338-satisfiesthermodynamiclaws, def:physics:phyx-mini-0338:phyxminiproblems-problemphyxmini0338-recordedanswerchoice, def:physics:phyx-mini-0338:phyxminiproblems-problemphyxmini0338-iscorrectdisplayedchoice}
The assigned autoformalize agent should translate this physics problem into Lean declarations in the covered file, with theorem and lemma proof bodies written as `by sorry`.
\end{theorem}
\begin{proof}
This is an autoformalization task, not a proof task. Produce statements that can later be proved by the physics prover without weakening the physical model.
\end{proof}

% --- Archon declaration topology begin ---
\section{Declaration topology}
\begin{definition}[Volume Quantity]
  \label{def:physics:phyx-mini-0338:phyxminiproblems-problemphyxmini0338-volumequantity}
  \lean{PhyXMiniProblems.ProblemPhyXMini0338.VolumeQuantity}
  A nonnegative physical volume, with dimension L³
\end{definition}
\begin{proof}
  This is the declared model definition of Volume Quantity.
\end{proof}
\begin{definition}[volume In Cubic Meters]
  \label{def:physics:phyx-mini-0338:phyxminiproblems-problemphyxmini0338-volumeincubicmeters}
  \lean{PhyXMiniProblems.ProblemPhyXMini0338.volumeInCubicMeters}
  \uses{def:physics:phyx-mini-0338:phyxminiproblems-problemphyxmini0338-volumequantity}
  Cubic-metre readout of a physical volume
\end{definition}
\begin{proof}
  This is the declared model definition of volume In Cubic Meters.
\end{proof}
\begin{definition}[pressure In Pascals]
  \label{def:physics:phyx-mini-0338:phyxminiproblems-problemphyxmini0338-pressureinpascals}
  \lean{PhyXMiniProblems.ProblemPhyXMini0338.pressureInPascals}
  Pascal readout of a physical pressure
\end{definition}
\begin{proof}
  This is the declared model definition of pressure In Pascals.
\end{proof}
\begin{definition}[energy In Joules]
  \label{def:physics:phyx-mini-0338:phyxminiproblems-problemphyxmini0338-energyinjoules}
  \lean{PhyXMiniProblems.ProblemPhyXMini0338.energyInJoules}
  Joule readout of a physical energy
\end{definition}
\begin{proof}
  This is the declared model definition of energy In Joules.
\end{proof}
\begin{definition}[Gas Model]
  \label{def:physics:phyx-mini-0338:phyxminiproblems-problemphyxmini0338-gasmodel}
  \lean{PhyXMiniProblems.ProblemPhyXMini0338.GasModel}
  Whether the sample obeys the ideal-gas model stated in the problem
\end{definition}
\begin{proof}
  This is the declared model definition of Gas Model.
\end{proof}
\begin{definition}[Gas Sample]
  \label{def:physics:phyx-mini-0338:phyxminiproblems-problemphyxmini0338-gassample}
  \lean{PhyXMiniProblems.ProblemPhyXMini0338.GasSample}
  \uses{def:physics:phyx-mini-0338:phyxminiproblems-problemphyxmini0338-gasmodel}
  A gas sample together with its measured amount-of-substance readout
\end{definition}
\begin{proof}
  This is the declared model definition of Gas Sample.
\end{proof}
\begin{definition}[State Label]
  \label{def:physics:phyx-mini-0338:phyxminiproblems-problemphyxmini0338-statelabel}
  \lean{PhyXMiniProblems.ProblemPhyXMini0338.StateLabel}
  The three labeled equilibrium states in the diagram
\end{definition}
\begin{proof}
  This is the declared model definition of State Label.
\end{proof}
\begin{definition}[Leg Label]
  \label{def:physics:phyx-mini-0338:phyxminiproblems-problemphyxmini0338-leglabel}
  \lean{PhyXMiniProblems.ProblemPhyXMini0338.LegLabel}
  The two directed legs of the process abc
\end{definition}
\begin{proof}
  This is the declared model definition of Leg Label.
\end{proof}
\begin{definition}[initial State]
  \label{def:physics:phyx-mini-0338:phyxminiproblems-problemphyxmini0338-leglabel-initialstate}
  \lean{PhyXMiniProblems.ProblemPhyXMini0338.LegLabel.initialState}
  \uses{def:physics:phyx-mini-0338:phyxminiproblems-problemphyxmini0338-statelabel, def:physics:phyx-mini-0338:phyxminiproblems-problemphyxmini0338-leglabel}
  Initial labeled state of each process leg
\end{definition}
\begin{proof}
  This is the declared model definition of initial State.
\end{proof}
\begin{definition}[final State]
  \label{def:physics:phyx-mini-0338:phyxminiproblems-problemphyxmini0338-leglabel-finalstate}
  \lean{PhyXMiniProblems.ProblemPhyXMini0338.LegLabel.finalState}
  \uses{def:physics:phyx-mini-0338:phyxminiproblems-problemphyxmini0338-statelabel, def:physics:phyx-mini-0338:phyxminiproblems-problemphyxmini0338-leglabel}
  Final labeled state of each process leg
\end{definition}
\begin{proof}
  This is the declared model definition of final State.
\end{proof}
\begin{definition}[PVState]
  \label{def:physics:phyx-mini-0338:phyxminiproblems-problemphyxmini0338-pvstate}
  \lean{PhyXMiniProblems.ProblemPhyXMini0338.PVState}
  \uses{def:physics:phyx-mini-0338:phyxminiproblems-problemphyxmini0338-volumequantity}
  A thermodynamic state as dimensionful pressure and volume
\end{definition}
\begin{proof}
  This is the declared model definition of PVState.
\end{proof}
\begin{definition}[Axis Quantity]
  \label{def:physics:phyx-mini-0338:phyxminiproblems-problemphyxmini0338-axisquantity}
  \lean{PhyXMiniProblems.ProblemPhyXMini0338.AxisQuantity}
  Physical quantity assigned to an axis of the supplied plot
\end{definition}
\begin{proof}
  This is the declared model definition of Axis Quantity.
\end{proof}
\begin{definition}[Volume Display Unit]
  \label{def:physics:phyx-mini-0338:phyxminiproblems-problemphyxmini0338-volumedisplayunit}
  \lean{PhyXMiniProblems.ProblemPhyXMini0338.VolumeDisplayUnit}
  Unit printed on the horizontal axis
\end{definition}
\begin{proof}
  This is the declared model definition of Volume Display Unit.
\end{proof}
\begin{definition}[Pressure Display Scale]
  \label{def:physics:phyx-mini-0338:phyxminiproblems-problemphyxmini0338-pressuredisplayscale}
  \lean{PhyXMiniProblems.ProblemPhyXMini0338.PressureDisplayScale}
  Scale printed on the vertical axis
\end{definition}
\begin{proof}
  This is the declared model definition of Pressure Display Scale.
\end{proof}
\begin{definition}[Leg Geometry]
  \label{def:physics:phyx-mini-0338:phyxminiproblems-problemphyxmini0338-leggeometry}
  \lean{PhyXMiniProblems.ProblemPhyXMini0338.LegGeometry}
  Geometric form of a directed leg in the pressure-volume plane
\end{definition}
\begin{proof}
  This is the declared model definition of Leg Geometry.
\end{proof}
\begin{definition}[PVDiagram]
  \label{def:physics:phyx-mini-0338:phyxminiproblems-problemphyxmini0338-pvdiagram}
  \lean{PhyXMiniProblems.ProblemPhyXMini0338.PVDiagram}
  \uses{def:physics:phyx-mini-0338:phyxminiproblems-problemphyxmini0338-statelabel, def:physics:phyx-mini-0338:phyxminiproblems-problemphyxmini0338-leglabel, def:physics:phyx-mini-0338:phyxminiproblems-problemphyxmini0338-pvstate, def:physics:phyx-mini-0338:phyxminiproblems-problemphyxmini0338-axisquantity, def:physics:phyx-mini-0338:phyxminiproblems-problemphyxmini0338-volumedisplayunit, def:physics:phyx-mini-0338:phyxminiproblems-problemphyxmini0338-pressuredisplayscale, def:physics:phyx-mini-0338:phyxminiproblems-problemphyxmini0338-leggeometry}
  The labeled pressure-volume diagram supplied with the problem
\end{definition}
\begin{proof}
  This is the declared model definition of PVDiagram.
\end{proof}
\begin{definition}[Ideal Gas Process ABC]
  \label{def:physics:phyx-mini-0338:phyxminiproblems-problemphyxmini0338-idealgasprocessabc}
  \lean{PhyXMiniProblems.ProblemPhyXMini0338.IdealGasProcessABC}
  \uses{def:physics:phyx-mini-0338:phyxminiproblems-problemphyxmini0338-gassample, def:physics:phyx-mini-0338:phyxminiproblems-problemphyxmini0338-leglabel, def:physics:phyx-mini-0338:phyxminiproblems-problemphyxmini0338-pvdiagram}
  The gas process and its energy transfers
\end{definition}
\begin{proof}
  This is the declared model definition of Ideal Gas Process ABC.
\end{proof}
\begin{definition}[Matches Problem And Primary Figure]
  \label{def:physics:phyx-mini-0338:phyxminiproblems-problemphyxmini0338-matchesproblemandprimaryfigure}
  \lean{PhyXMiniProblems.ProblemPhyXMini0338.MatchesProblemAndPrimaryFigure}
  \uses{def:physics:phyx-mini-0338:phyxminiproblems-problemphyxmini0338-volumeincubicmeters, def:physics:phyx-mini-0338:phyxminiproblems-problemphyxmini0338-pressureinpascals, def:physics:phyx-mini-0338:phyxminiproblems-problemphyxmini0338-energyinjoules, def:physics:phyx-mini-0338:phyxminiproblems-problemphyxmini0338-leglabel, def:physics:phyx-mini-0338:phyxminiproblems-problemphyxmini0338-idealgasprocessabc}
  Problem data and exact grid readouts from the primary image. The amount is 0.450 mol, and the specified internal-energy increase is 15000 J. No value of the heat added or of the pressure-volume work is assumed here
\end{definition}
\begin{proof}
  This is the declared model definition of Matches Problem And Primary Figure.
\end{proof}
\begin{definition}[Has Physical Process Parameters]
  \label{def:physics:phyx-mini-0338:phyxminiproblems-problemphyxmini0338-hasphysicalprocessparameters}
  \lean{PhyXMiniProblems.ProblemPhyXMini0338.HasPhysicalProcessParameters}
  \uses{def:physics:phyx-mini-0338:phyxminiproblems-problemphyxmini0338-volumeincubicmeters, def:physics:phyx-mini-0338:phyxminiproblems-problemphyxmini0338-pressureinpascals, def:physics:phyx-mini-0338:phyxminiproblems-problemphyxmini0338-statelabel, def:physics:phyx-mini-0338:phyxminiproblems-problemphyxmini0338-idealgasprocessabc}
  Positivity and ordering conditions for the physical process
\end{definition}
\begin{proof}
  This is the declared model definition of Has Physical Process Parameters.
\end{proof}
\begin{definition}[Satisfies Thermodynamic Laws]
  \label{def:physics:phyx-mini-0338:phyxminiproblems-problemphyxmini0338-satisfiesthermodynamiclaws}
  \lean{PhyXMiniProblems.ProblemPhyXMini0338.SatisfiesThermodynamicLaws}
  \uses{def:physics:phyx-mini-0338:phyxminiproblems-problemphyxmini0338-volumeincubicmeters, def:physics:phyx-mini-0338:phyxminiproblems-problemphyxmini0338-pressureinpascals, def:physics:phyx-mini-0338:phyxminiproblems-problemphyxmini0338-energyinjoules, def:physics:phyx-mini-0338:phyxminiproblems-problemphyxmini0338-leglabel, def:physics:phyx-mini-0338:phyxminiproblems-problemphyxmini0338-idealgasprocessabc}
  The pressure-volume work law and the first law of thermodynamics. For a straight leg, the integral ∫ p dV is the average endpoint pressure times the volume change. A vertical isochoric leg performs zero work. The final field uses the sign convention Q = ΔU + W\_by\_gas. All equations are symbolic governing laws; none fixes the requested heat to a numerical answer
\end{definition}
\begin{proof}
  This is the declared model definition of Satisfies Thermodynamic Laws.
\end{proof}
\begin{lemma}[work Done By Gas joules eq 21000]
  \label{lem:physics:phyx-mini-0338:phyxminiproblems-problemphyxmini0338-workdonebygas-joules-eq-21000}
  \lean{PhyXMiniProblems.ProblemPhyXMini0338.workDoneByGas_joules_eq_21000}
  \uses{def:physics:phyx-mini-0338:phyxminiproblems-problemphyxmini0338-energyinjoules, def:physics:phyx-mini-0338:phyxminiproblems-problemphyxmini0338-idealgasprocessabc, def:physics:phyx-mini-0338:phyxminiproblems-problemphyxmini0338-matchesproblemandprimaryfigure, def:physics:phyx-mini-0338:phyxminiproblems-problemphyxmini0338-hasphysicalprocessparameters, def:physics:phyx-mini-0338:phyxminiproblems-problemphyxmini0338-satisfiesthermodynamiclaws}
  The net work read from the straight ab leg and the isochoric bc leg is 21000 J. This is a derived intermediate result, not a figure readout
\end{lemma}
\begin{proof}
  The conclusion follows from the governing relations and figure readouts named in the statement of work Done By Gas joules eq 21000.
\end{proof}
\begin{definition}[Answer Choice]
  \label{def:physics:phyx-mini-0338:phyxminiproblems-problemphyxmini0338-answerchoice}
  \lean{PhyXMiniProblems.ProblemPhyXMini0338.AnswerChoice}
  Labels printed beside the four candidate heat values
\end{definition}
\begin{proof}
  This is the declared model definition of Answer Choice.
\end{proof}
\begin{definition}[heat In Joules]
  \label{def:physics:phyx-mini-0338:phyxminiproblems-problemphyxmini0338-answerchoice-heatinjoules}
  \lean{PhyXMiniProblems.ProblemPhyXMini0338.AnswerChoice.heatInJoules}
  \uses{def:physics:phyx-mini-0338:phyxminiproblems-problemphyxmini0338-answerchoice}
  Heat in joules printed beside each answer label
\end{definition}
\begin{proof}
  This is the declared model definition of heat In Joules.
\end{proof}
\begin{definition}[recorded Answer Choice]
  \label{def:physics:phyx-mini-0338:phyxminiproblems-problemphyxmini0338-recordedanswerchoice}
  \lean{PhyXMiniProblems.ProblemPhyXMini0338.recordedAnswerChoice}
  \uses{def:physics:phyx-mini-0338:phyxminiproblems-problemphyxmini0338-answerchoice}
  The answer label recorded in the supplied dataset metadata
\end{definition}
\begin{proof}
  This is the declared model definition of recorded Answer Choice.
\end{proof}
\begin{definition}[Is Correct Displayed Choice]
  \label{def:physics:phyx-mini-0338:phyxminiproblems-problemphyxmini0338-iscorrectdisplayedchoice}
  \lean{PhyXMiniProblems.ProblemPhyXMini0338.IsCorrectDisplayedChoice}
  \uses{def:physics:phyx-mini-0338:phyxminiproblems-problemphyxmini0338-energyinjoules, def:physics:phyx-mini-0338:phyxminiproblems-problemphyxmini0338-idealgasprocessabc, def:physics:phyx-mini-0338:phyxminiproblems-problemphyxmini0338-answerchoice}
  A displayed choice agrees with the physical heat added in the process
\end{definition}
\begin{proof}
  This is the declared model definition of Is Correct Displayed Choice.
\end{proof}
\begin{theorem}[heat Added During Process ABC joules eq 36000]
  \label{thm:physics:phyx-mini-0338:phyxminiproblems-problemphyxmini0338-heataddedduringprocessabc-joules-eq-36000}
  \lean{PhyXMiniProblems.ProblemPhyXMini0338.heatAddedDuringProcessABC_joules_eq_36000}
  \uses{def:physics:phyx-mini-0338:phyxminiproblems-problemphyxmini0338-energyinjoules, def:physics:phyx-mini-0338:phyxminiproblems-problemphyxmini0338-idealgasprocessabc, def:physics:phyx-mini-0338:phyxminiproblems-problemphyxmini0338-matchesproblemandprimaryfigure, def:physics:phyx-mini-0338:phyxminiproblems-problemphyxmini0338-hasphysicalprocessparameters, def:physics:phyx-mini-0338:phyxminiproblems-problemphyxmini0338-satisfiesthermodynamiclaws}
  The heat that must be added during a → b → c is 36000 J. This is the declaration corresponding to
\end{theorem}
\begin{proof}
  The conclusion follows from the governing relations and figure readouts named in the statement of heat Added During Process ABC joules eq 36000.
\end{proof}
% --- Archon declaration topology end ---
% --- Archon physics formalization source end ---
